\section{Tercera vía de extensión: Inference-Time Scaling}

\subsection{Con los parámetros fijos, aún se puede cambiar más cómputo por tasa de éxito}

El punto de partida de Large Language Monkeys es muy directo: muestrear la misma pregunta de forma independiente muchas veces; mientras el modelo tenga una probabilidad de éxito distinta de cero para esa pregunta, es posible acertar la solución correcta en alguna muestra. El sistema se divide en dos pasos:

\begin{enumerate}
    \item \textbf{Generation}: generar una gran cantidad de candidatos, buscando coverage;
    \item \textbf{Verification}: identificar la solución correcta entre los candidatos, buscando precision.
\end{enumerate}

\begin{figure}[H]
\centering
\includegraphics[width=0.96\textwidth]{figures/language-monkeys.jpg}
\caption{Large Language Monkeys: generación diversa y selección mediante verificación.\protect\footnotemark}
\end{figure}
\footnotetext{Intervalo de la explicación en video: 00:20:03--00:22:43.}

Supongamos que la probabilidad de éxito independiente de una muestra es $p$; la probabilidad de tener al menos un éxito al muestrear $k$ veces es:

$$
P_{\mathrm{solve}}(k)=1-(1-p)^k.
$$

\begin{itemize}
    \item $p$: probabilidad de obtener un candidato correcto en una sola muestra;
    \item $k$: presupuesto de muestreo en el momento de la inferencia;
    \item $(1-p)^k$: probabilidad de fallar las $k$ veces;
    \item $P_{\mathrm{solve}}(k)$: probabilidad de que exista al menos una solución correcta en el conjunto de candidatos.
\end{itemize}

Cuando $p$ es muy pequeño, se necesita un $k$ muy grande para ver una ganancia apreciable; cuando $p=0$, muestrear más no sirve de nada. Esto explica el límite del test-time scaling: puede explotar capacidades que el modelo \textbf{ya posee con baja probabilidad}, pero no puede crear de la nada un algoritmo que el modelo no sabe en absoluto.

\begin{importantbox}{Coverage y precisión final no son lo mismo}
Oracle pass@$k$ supone que podemos saber cuál de los candidatos es correcto, y mide la cobertura de la generación; un sistema real todavía enfrenta el problema de la selección. Si el verifier no logra identificar la respuesta correcta, un conjunto de candidatos por bueno que sea no se traduce en el pass@1 final. Esta diferencia es precisamente el generation--verification gap.
\end{importantbox}

\subsection{Por qué incluso un modelo pequeño puede ser ``amplificado'' con muestreo masivo}

Los experimentos muestran que, para distintos modelos y distintas preguntas, a menudo existe una probabilidad de éxito de cola larga: un modelo de 8B puede fallar la mayoría de las veces en una pregunta, pero acertar en trayectorias muy escasas. Aumentar el número de muestras va recolectando poco a poco esas trayectorias de éxito poco frecuentes, de modo que, bajo una evaluación oracle, el modelo pequeño se acerca o incluso supera el desempeño de un solo intento de un modelo más fuerte.

\begin{figure}[H]
\centering
\includegraphics[width=0.97\textwidth]{figures/repeated-sampling-results.jpg}
\caption{El muestreo repetido amplía la cobertura de candidatos en tareas de matemáticas, código y demostraciones formales.\protect\footnotemark}
\end{figure}
\footnotetext{Intervalo de la explicación en video: 00:23:49--00:27:12.}

Pero la ganancia de cómputo tiene rendimientos decrecientes: una vez cubiertas las preguntas fáciles, el muestreo adicional se consume sobre todo en las preguntas extremadamente difíciles. El diseño del sistema no debe preguntarse solo ``¿muestrear más mejora el resultado?'', sino también ``¿cuánta tasa de éxito marginal se obtiene por cada unidad de token, latencia y energía consumidos?''.

\subsection{El cómputo de entrenamiento y el cómputo de prueba son dos presupuestos distintos}

Los reasoning model incorporan a la inferencia cadenas de pensamiento más largas, búsqueda y autocorrección. El curso usa las curvas de entrenamiento y de tiempo de prueba de o1 para ilustrar que la capacidad puede crecer tanto con el cómputo de entrenamiento como, una vez fijado el modelo, con el test-time compute.

\begin{figure}[H]
\centering
\includegraphics[width=0.9\textwidth]{figures/o1-scaling.jpg}
\caption{Tanto la fase de entrenamiento como la fase de prueba tienen un eje de cómputo escalable.\protect\footnotemark}
\end{figure}
\footnotetext{Intervalo de la explicación en video: 00:29:40--00:32:07.}

\begin{figure}[H]
\centering
\includegraphics[width=0.94\textwidth]{figures/reasoning-models.jpg}
\caption{Los reasoning model consumen más cómputo intermedio mediante descomposición, evaluación, corrección y alternativas.\protect\footnotemark}
\end{figure}
\footnotetext{Intervalo de la explicación en video: 00:32:08--00:33:55.}

\begin{knowledgebox}{Escalamiento paralelo y escalamiento en serie}
El \textbf{escalamiento paralelo} genera de forma independiente varios candidatos; es fácil de paralelizar, pero los candidatos no comparten experiencia entre sí. El \textbf{escalamiento en serie} hace que los pasos posteriores lean los resultados anteriores y realicen revisión, búsqueda o retroceso; puede ser más eficiente, pero si el estado inicial es erróneo, el sesgo también se propaga a los pasos siguientes. Los sistemas reales suelen combinar ambos.
\end{knowledgebox}

\subsection{El presupuesto de inferencia debe adaptarse a la dificultad de la pregunta}

Asignar el mismo presupuesto de tokens a cada pregunta desperdicia cómputo: las preguntas sencillas se resuelven en un solo intento, mientras que las difíciles requieren más exploración. Un controlador más razonable debería decidir si continuar el cómputo con base en la incertidumbre, la puntuación de verificación, la consistencia entre candidatos o señales intermedias de fallo.

El presupuesto de asignación puede escribirse como una optimización con restricciones:

$$
\max_{k_1,\ldots,k_n}\sum_{i=1}^{n} U_i(k_i)
\quad\text{s.t.}\quad
\sum_{i=1}^{n} c_i k_i\le B.
$$

\begin{itemize}
    \item $k_i$: número de pasos de inferencia o de muestras que recibe la pregunta $i$;
    \item $U_i(k_i)$: utilidad esperada que aporta el cómputo adicional;
    \item $c_i$: costo de cada paso;
    \item $B$: presupuesto total de tokens, latencia o dinero.
\end{itemize}

\subsection{Resumen de la sección}

Inference-time scaling traslada el ``más cómputo'' del clúster de entrenamiento a cada solicitud individual. El repeated sampling mejora la coverage, el reasoning/search mejora la calidad de cada trayectoria individual, y el verifier determina si estos candidatos adicionales pueden traducirse en la precisión final. Muchos de los métodos que se verán a lo largo del curso pueden entenderse como mejoras a estas tres partes.
